\section{审判管辖}
\begin{enumerate}
  \item 用排除法。
  \item 级别管辖。默认基层。
        \begin{enumerate}
          \item 中院。\bs{危恐无死没缺}。危害国家安全、恐怖活动犯罪、可能判无期死刑、没收违法所得程序、缺席审判。
          \item 高院。省内重大。
          \item 最高法。全国重大。
          \item 级别变通。上可以审下，下不可审上。\bs{数罪或数人，就高不就低}。\bs{就高例外。未成年分案处理}。
        \end{enumerate}
  \item 地域管辖。
        \begin{enumerate}
          \item 犯罪地（行为地 + 结果地）为主，被告人居住地为辅。
          \item 最初受理地为主，主要犯罪地为辅。
          \item 要下结论 \yjt \zd{共同管辖，先到先得}。
        \end{enumerate}
  \item 专门管辖。一般了解即可。
  \item 指定管辖。\al{院长涉及回避、请求全员/整体回避（当前于法无据）}。
        \begin{enumerate}
          \item 管辖不明、不宜、规避管辖。协商不成，分别层报共同上级来指定管辖。\al{甲宣称 A 市中级法院院长乙为当地黑社会“保护伞”}。
          \item 案件材料移送。公诉要退回，自诉直接送。
        \end{enumerate}
  \item \bs{问管辖有无问题}。\zd{立案管辖、审判管辖（级别、地域、有无指定管辖）、特殊案件管辖、是否回避}都要答到。
\end{enumerate}
